% ---
% id: reductions/rem:notchannel
% kind: rem
% title: "Reductions that are not squashings"
% summary: >-
%   The universal squash model establishes a statistical equivalence and
%   cutoff-based dimension reduction controls the effect of a cutoff, neither
%   by exhibiting a squashing channel. Without such a channel the output-state
%   equality behind [[thm:transfer]] is not available.
% prominence: standard
% status: accepted
% uses: []
% mentions: [thm:transfer, thm:factorization, req:channel-witness, req:certified-entries-only]
% assumes: []
% refs:
%   - "fung2011universal, Sec. II"
%   - "upadhyaya2021dimension, Sec. III"
%   - "nahar2024postselection, Lemma 6"
% provenance: newly-derived
% lean:
% ---

\begin{remark}\label{rem:notchannel}
Not every dimension reduction used in QKD comes with a squashing channel, and
one that does not cannot be fed into the security transfer,
Theorem~\ref{thm:transfer}. The universal squash model of
\cite[Sec.~II]{fung2011universal} establishes a statistical
equivalence between the threshold-detector setup and an ideal squashed one by
routing both through classical post-processing of a photon-number-resolving
measurement, and the dimension-reduction method of
\cite[Sec.~III]{upadhyaya2021dimension} controls the effect of a finite cutoff
rather than exhibiting a channel. Neither proceeds through the existence of a
$\Lambda$ with $\Lambda^{\dagger}(F_v)=M_v$, which is exactly the hypothesis the
output-state equality of Theorem~\ref{thm:factorization} runs on and the
hypothesis \cite[Lemma~6]{nahar2024postselection} states for a narrower model.
The rule that follows for the emitted artefact, a \tref{reduction chain} of
links, is Requirement~\ref{req:channel-witness}: an explicit
\tref{channel witness} for every link. The companion rule for constraint
entries, exactly one named certificate each, is
Requirement~\ref{req:certified-entries-only}.
\end{remark}
